\documentclass[a4paper,12pt]{article}

% =======================
% IMPORTATION DES PACKAGE
% =======================

\usepackage[utf8]{inputenc}
\usepackage[T1]{fontenc}
\usepackage[margin=2cm]{geometry}
\usepackage{amsfonts, amsmath, amssymb}
\usepackage{lmodern}
\usepackage{theorem}
\usepackage[french]{babel}

\pagestyle{headings}

\date{\today}
\author{}
\title{RÉSOLUTION D'ÉQUATION NON LINÉAIRE}

\renewcommand{\familydefault}{\sfdefault}

\theoremstyle{break}
\theorembodyfont{\upshape}
\theoremheaderfont{\scshape}

\newtheorem{meth}{Méthode}


% ==================
% DOCUMENT PRICIPALE
% ==================

\begin{document}
\maketitle

Le but est de décrire des algos pour résoudre des équation non liniéaire de type $f(x)=0$.
Dans cette depot, je vous presente des algorithmes pour resoudre des equations non lineaire.
Dans mon depot, il y en a 3 mais il y en a d'autre mais pour l'instant vous puvez vous contenter de ces 3 algorithmes suivant:
\begin{itemize}
    \item Méthode de Dichotomie,
    \item Méthode de Newton,
    \item Méthode du point fixe.
\end{itemize}

\begin{meth}[Dichotomie]
    \textsc{Algorithme:} \\
        Soit $f : [a_{0}, b_{0}] \rightarrow \mathbb{R}$ une fonction continue monotone telle que $f(a_{0})f(b_{0}) < 0$. \\
        Pour $m = 1, \cdots, N$ faire: \\ $m = \frac{a_{n} + b_{n}}{2}$
        \begin{itemize}
            \item $f(a_{n})f(m)<0, a_{n+1} = a_{n} \mbox{ et } b_{n} = m$,
            \item Sinon $a_{n+1} = m \mbox{ et } b_{n+1} = b_{n}$.
        \end{itemize}
        On a : $a_{n+1} - b_{n+1} = \frac{a_{n} - b_{n}}{2}$. \\
        Soit $a_{n} - b_{n} = \frac{a_{0} - b_{0}}{2^{n}} \rightarrow 0$. \\
        On peut choisir un temps d'arrêt telle que: $\frac{a_{0} - b_{0}}{2^{N}} < \epsilon$.
\end{meth}

\begin{meth}[Newton]
On assimille à une tangente en un point $(x_{n}, f(x_{n}))$, soit la droite d'équation $Y = f(x_{0}) + f^{'}(x_{n})(x - x_{n})$. \\
    \textsc{Algorithme:} \\
    Pour $n = 0,1,2, \cdots$, on a $$x_{n+1} = x_{n} - \frac{f(x_{n})}{f^{'}(x_{n})}.$$
\end{meth}

\begin{meth}[Point fixe]
    Elle consiste d'abord à remplacer $f(x)=0$ par une fonction $g(x)=x$ ayant même solution.
    \textsc{Algorithme:} \\
        Pour $n=1,2,3, \cdots$, on a: $$x_{n+1} = g(x_{n}).$$
\end{meth}
\end{document}